\section{Компьютерная графика. Работа с графическими библиотеками}

Компьютерная графика --- область информатики, изучающая методы создания,
хранения, обработки и визуального представления изображений с помощью вычислительной
техники.

В общем случае графическая система преобразует некоторое математическое или
информационное описание сцены в изображение:
\[
\text{модель сцены}
\longrightarrow
\text{геометрические преобразования}
\longrightarrow
\text{проекция}
\longrightarrow
\text{растеризация}
\longrightarrow
\text{изображение}.
\]

Компьютерная графика используется не только для получения реалистичных изображений,
но и для визуального анализа результатов вычислений, математического моделирования,
САПР, научной визуализации и создания пользовательских интерфейсов.

\subsection{Основные задачи компьютерной графики}

К основным задачам компьютерной графики относятся:

\begin{itemize}
	\item \textbf{синтез изображений} --- построение изображения по математическому
	описанию сцены;
	
	\item \textbf{геометрическое моделирование} --- описание формы и положения
	объектов в двумерном или трёхмерном пространстве;
	
	\item \textbf{визуализация данных} --- представление численных данных в форме,
	удобной для анализа человеком;
	
	\item \textbf{обработка изображений} --- изменение уже существующих изображений:
	фильтрация, масштабирование, коррекция, выделение границ и т.д.;
	
	\item \textbf{анимация} --- формирование последовательности изображений,
	описывающей изменение сцены во времени;
	
	\item \textbf{интерактивная графика} --- изменение изображения в ответ на действия
	пользователя;
	
	\item \textbf{фотореалистичный рендеринг} --- моделирование распространения света
	для получения изображения, близкого к реальному.
\end{itemize}

Следует различать компьютерную графику и компьютерное зрение. В компьютерной
графике обычно решается задача
\[
\text{описание сцены} \rightarrow \text{изображение},
\]
а в компьютерном зрении --- в некотором смысле обратная задача:
\[
\text{изображение} \rightarrow \text{информация о сцене}.
\]

Для математического моделирования особенно важна \textbf{научная визуализация}:
построение графиков функций, полей, поверхностей, сеток, линий уровня, траекторий,
изоповерхностей и других представлений результатов численного расчёта.

\subsection{Растровая и векторная графика}

Существуют два основных способа представления графической информации:
растровый и векторный.

\subsubsection*{Растровая графика}

Растровое изображение представляет собой прямоугольную матрицу элементов ---
\textbf{пикселей}:
\[
I = \{c_{ij}\},
\]
где $c_{ij}$ задаёт цвет пикселя с координатами $(i,j)$.

Если используется цветовая модель RGB, то пиксель можно представить тройкой
\[
c_{ij}=(R_{ij},G_{ij},B_{ij}).
\]

Часто добавляется четвёртая компонента $\alpha$, задающая прозрачность:
\[
c_{ij}=(R,G,B,\alpha).
\]

Количество битов, используемых для хранения цвета одного пикселя, называется
\textbf{глубиной цвета}. Например, при использовании $8$ бит на каждый канал RGB
получаем
\[
2^{24}
\]
возможных цветов.

Основные характеристики растрового изображения:

\begin{itemize}
	\item ширина и высота в пикселях;
	\item разрешение;
	\item глубина цвета;
	\item цветовая модель;
	\item наличие альфа-канала.
\end{itemize}

Преимуществами растровой графики являются простота представления и возможность
описывать сложные изображения с большим количеством мелких деталей.

Недостатком является зависимость от разрешения: при значительном увеличении
изображения становятся заметны отдельные пиксели.

\subsubsection*{Векторная графика}

Векторное изображение задаётся набором геометрических примитивов:

\begin{itemize}
	\item точек;
	\item отрезков;
	\item кривых;
	\item многоугольников;
	\item окружностей;
	\item текста.
\end{itemize}

Например, отрезок можно задать двумя точками
\[
A=(x_1,y_1), \qquad B=(x_2,y_2),
\]
а окружность --- центром $(x_0,y_0)$ и радиусом $r$:
\[
(x-x_0)^2+(y-y_0)^2=r^2.
\]

Векторная графика не зависит непосредственно от разрешения устройства вывода.
При изменении масштаба геометрические примитивы строятся заново.

Преимущества векторной графики:

\begin{itemize}
	\item качественное масштабирование;
	\item компактное представление простых геометрических объектов;
	\item удобство редактирования отдельных элементов.
\end{itemize}

Недостатком является сложность описания фотографий и других изображений с большим
количеством нерегулярных деталей.

При отображении на экран векторное изображение в конечном итоге также должно быть
преобразовано в набор пикселей. Этот процесс называется \textbf{растеризацией}.

\subsection{Геометрические преобразования и системы координат}

Основными геометрическими преобразованиями являются:

\begin{itemize}
	\item перенос;
	\item масштабирование;
	\item вращение;
	\item отражение;
	\item проецирование.
\end{itemize}

Для удобного представления преобразований используются
\textbf{однородные координаты}.

Точка трёхмерного пространства
\[
(x,y,z)
\]
представляется четырёхмерным вектором
\[
\mathbf{p}=
\begin{pmatrix}
	x\\y\\z\\1
\end{pmatrix}.
\]

Тогда любое аффинное преобразование можно записать как умножение на матрицу:
\[
\mathbf{p}'=M\mathbf{p}.
\]

Матрица переноса имеет вид
\[
T=
\begin{pmatrix}
	1&0&0&t_x\\
	0&1&0&t_y\\
	0&0&1&t_z\\
	0&0&0&1
\end{pmatrix}.
\]

Матрица масштабирования:
\[
S=
\begin{pmatrix}
	s_x&0&0&0\\
	0&s_y&0&0\\
	0&0&s_z&0\\
	0&0&0&1
\end{pmatrix}.
\]

Например, вращение вокруг оси $z$ на угол $\varphi$ задаётся матрицей
\[
R_z=
\begin{pmatrix}
	\cos\varphi&-\sin\varphi&0&0\\
	\sin\varphi& \cos\varphi&0&0\\
	0&0&1&0\\
	0&0&0&1
\end{pmatrix}.
\]

Несколько преобразований можно объединить:
\[
M=M_3M_2M_1.
\]

При этом порядок умножения матриц существенен, поскольку в общем случае
\[
M_1M_2\ne M_2M_1.
\]

В трёхмерной графике обычно последовательно используются несколько систем
координат:

\begin{enumerate}
	\item \textbf{локальные координаты объекта} --- координаты относительно самого
	объекта;
	
	\item \textbf{мировые координаты} --- положение объекта в общей сцене;
	
	\item \textbf{координаты камеры} --- система координат наблюдателя;
	
	\item \textbf{координаты отсечения} (\textit{clip coordinates});
	
	\item \textbf{нормализованные координаты устройства} (\textit{NDC});
	
	\item \textbf{экранные координаты}.
\end{enumerate}

Типичная последовательность преобразований вершины имеет вид
\[
p_{\mathrm{clip}}
=
PVMp_{\mathrm{local}},
\]
где

\begin{itemize}
	\item $M$ --- матрица модели;
	\item $V$ --- матрица вида;
	\item $P$ --- матрица проекции.
\end{itemize}

Для трёхмерных сцен применяются два основных вида проекции.

При \textbf{ортографической проекции} размер объекта практически не зависит от
расстояния до камеры.

При \textbf{перспективной проекции} удалённые объекты выглядят меньше.
После перспективного преобразования выполняется деление на однородную координату:
\[
x_{\mathrm{NDC}}=\frac{x_c}{w_c},
\qquad
y_{\mathrm{NDC}}=\frac{y_c}{w_c},
\qquad
z_{\mathrm{NDC}}=\frac{z_c}{w_c}.
\]

\subsection{Графический конвейер}

\textbf{Графический конвейер} (\textit{graphics pipeline}) --- последовательность
этапов преобразования описания сцены в растровое изображение.

Для современного растеризационного графического API конвейер можно схематически
представить следующим образом:
\[
\begin{split}
	&\text{вершины}
	\rightarrow
	\text{обработка вершин}
	\rightarrow
	\text{сборка примитивов}
	\rightarrow\\
	&\rightarrow
	\text{отсечение}
	\rightarrow
	\text{растеризация}
	\rightarrow
	\text{обработка фрагментов}
	\rightarrow
	\text{кадр}.
\end{split}
\]

Исходная геометрия обычно представляется набором вершин. Вершина может содержать:

\begin{itemize}
	\item координаты;
	\item нормаль;
	\item цвет;
	\item текстурные координаты;
	\item дополнительные пользовательские атрибуты.
\end{itemize}

На стадии обработки вершин координаты преобразуются из локальной системы координат
в координаты отсечения.

Затем вершины собираются в графические примитивы:

\begin{itemize}
	\item точки;
	\item линии;
	\item треугольники.
\end{itemize}

Треугольник является основным примитивом трёхмерной растровой графики, поскольку:

\begin{itemize}
	\item три точки всегда определяют плоскость;
	\item треугольник всегда выпуклый;
	\item его легко растеризовать;
	\item произвольную полигональную поверхность можно разбить на треугольники.
\end{itemize}

Примитивы, полностью находящиеся вне области видимости, отбрасываются.
Пересекающие границу области видимости примитивы подвергаются
\textbf{отсечению}.

После этого производится переход к экранным координатам и растеризация.

\subsection{Растеризация и визуализация}

\textbf{Растеризация} --- преобразование геометрического примитива в множество
фрагментов, соответствующих потенциальным пикселям экрана.

Например, при растеризации треугольника требуется определить множество пикселей,
центры которых принадлежат его проекции.

Атрибуты вершин внутри треугольника интерполируются. Для этого удобно использовать
\textbf{барицентрические координаты}:
\[
p=\lambda_1p_1+\lambda_2p_2+\lambda_3p_3,
\]
где
\[
\lambda_1+\lambda_2+\lambda_3=1.
\]

Тогда некоторый атрибут $a$, заданный в вершинах, можно интерполировать:
\[
a(p)
=
\lambda_1a_1+\lambda_2a_2+\lambda_3a_3.
\]

На стадии обработки фрагментов обычно определяются:

\begin{itemize}
	\item цвет;
	\item прозрачность;
	\item значение глубины;
	\item параметры освещения;
	\item значения текстур.
\end{itemize}

Для удаления невидимых поверхностей используется
\textbf{буфер глубины} (\textit{Z-buffer}).

Для каждого пикселя хранится глубина ближайшего уже нарисованного фрагмента.
Новый фрагмент выводится только в случае успешного теста глубины, например
\[
z_{\mathrm{new}} < z_{\mathrm{buffer}}.
\]

После этого значение глубины обновляется.

Для вычисления цвета поверхности могут использоваться модели освещения.
Например, простая диффузная составляющая модели Ламберта имеет вид
\[
I_d=k_d I_L \max(0,\mathbf{n}\cdot\mathbf{l}),
\]
где

\begin{itemize}
	\item $\mathbf{n}$ --- единичная нормаль к поверхности;
	\item $\mathbf{l}$ --- направление на источник света;
	\item $I_L$ --- интенсивность источника;
	\item $k_d$ --- коэффициент диффузного отражения.
\end{itemize}

В реальных системах также применяются текстурирование, смешивание цветов
(\textit{blending}), сглаживание границ (\textit{antialiasing}), карты нормалей,
теневые карты и другие методы.

Следует отличать \textbf{растеризацию} от более общего понятия
\textbf{рендеринга}. Рендеринг --- процесс получения изображения сцены в целом,
а растеризация --- лишь один из возможных способов определить вклад геометрических
примитивов в пиксели.

Другим важным классом методов является \textbf{трассировка лучей}. В простейшем
варианте для каждого пикселя строится луч
\[
r(t)=o+td,
\]
где $o$ --- положение камеры, а $d$ --- направление луча. Затем определяется
ближайшее пересечение луча с объектами сцены.

Трассировка лучей естественным образом позволяет моделировать отражения,
преломления и тени, но обычно требует большего объёма вычислений, чем
растеризация.

\subsection{Графические программные интерфейсы}

Для работы с графическим оборудованием используются
\textbf{графические программные интерфейсы} --- API
(\textit{Application Programming Interface}).

К распространённым графическим API относятся:

\begin{itemize}
	\item OpenGL;
	\item Vulkan;
	\item Direct3D;
	\item Metal;
	\item WebGL;
	\item WebGPU.
\end{itemize}

Графический API предоставляет программе средства для:

\begin{itemize}
	\item создания графического контекста или устройства;
	\item передачи данных в память GPU;
	\item создания буферов вершин и индексов;
	\item загрузки текстур;
	\item создания и запуска шейдеров;
	\item настройки графического конвейера;
	\item формирования команд рисования;
	\item синхронизации CPU и GPU;
	\item вывода сформированного кадра на экран.
\end{itemize}

Например, геометрия может храниться в \textbf{вершинном буфере}
(\textit{vertex buffer}), а порядок соединения вершин в примитивы ---
в \textbf{индексном буфере}.

Если имеются вершины
\[
v_0,v_1,\ldots,v_{N-1},
\]
то треугольники можно задавать индексами:
\[
(0,1,2),\quad (2,3,0),\quad \ldots
\]

Это позволяет не хранить одну и ту же вершину несколько раз.

Графические API можно условно разделить на высокоуровневые и низкоуровневые.

OpenGL традиционно скрывает значительную часть управления ресурсами и
синхронизацией за внутренним состоянием драйвера.

Vulkan, Direct3D~12 и Metal предоставляют программисту значительно более явный
контроль над:

\begin{itemize}
	\item памятью;
	\item очередями команд;
	\item синхронизацией;
	\item состоянием графического конвейера.
\end{itemize}

Это позволяет уменьшить накладные расходы и лучше использовать современное
многоядерное оборудование, но усложняет программу.

Над низкоуровневыми API часто используются более высокоуровневые библиотеки и
движки. Например, для научной визуализации применяются библиотеки, которые
скрывают значительную часть работы с GPU и позволяют строить графики,
поверхности и поля непосредственно из численных данных.

\subsection{Шейдеры и программируемый конвейер}

Современные GPU используют \textbf{программируемый графический конвейер}.
Часть стадий обработки графики задаётся небольшими программами ---
\textbf{шейдерами}.

Основными типами шейдеров являются:

\begin{itemize}
	\item вершинный шейдер;
	\item фрагментный (пиксельный) шейдер;
	\item геометрический шейдер;
	\item тесселяционные шейдеры;
	\item вычислительный шейдер.
\end{itemize}

\textbf{Вершинный шейдер} выполняется для каждой вершины. Одна из его основных
задач --- вычислить положение вершины в координатах отсечения:
\[
p_{\mathrm{clip}}=PVMp.
\]

Кроме положения вершины он может преобразовывать нормали, вычислять текстурные
координаты и передавать параметры следующим стадиям конвейера.

\textbf{Фрагментный шейдер} выполняется для фрагментов, полученных после
растеризации. Обычно он вычисляет итоговый цвет:
\[
C = F(\text{normal},\text{light},\text{texture},\ldots).
\]

Например, значение цвета текстуры определяется по текстурным координатам
\[
(u,v).
\]

Изображение-текстура рассматривается как функция
\[
T(u,v),
\]
возвращающая цвет.

Шейдеры выполняются параллельно для большого количества вершин или фрагментов.
Именно массовый параллелизм является одной из причин высокой эффективности GPU
в задачах компьютерной графики.

\textbf{Вычислительные шейдеры} отличаются от обычных графических шейдеров тем,
что не обязаны непосредственно участвовать в рисовании примитивов. Они позволяют
использовать GPU для вычислений общего назначения, например:

\begin{itemize}
	\item обработки изображений;
	\item моделирования частиц;
	\item вычисления физических полей;
	\item подготовки данных для визуализации;
	\item некоторых численных алгоритмов.
\end{itemize}

Поэтому современная граница между графическими вычислениями и вычислениями общего
назначения на GPU является достаточно условной.

\subsection{Научная визуализация и интерактивная графика}

В математическом моделировании компьютерная графика является инструментом анализа
результатов вычислительного эксперимента.

В зависимости от размерности и природы данных применяются различные способы
визуализации.

Для скалярной функции одной переменной
\[
y=f(x)
\]
обычно строится график функции.

Для функции двух переменных
\[
z=f(x,y)
\]
можно использовать:

\begin{itemize}
	\item поверхность;
	\item карту цветов;
	\item линии уровня.
\end{itemize}

Линией уровня является множество точек
\[
f(x,y)=C.
\]

Для скалярного трёхмерного поля
\[
f(x,y,z)
\]
используются сечения, объёмная визуализация и \textbf{изоповерхности}
\[
f(x,y,z)=C.
\]

Для векторного поля
\[
\mathbf{v}(x,y,z)
\]
могут применяться:

\begin{itemize}
	\item стрелочные диаграммы;
	\item линии тока;
	\item траектории;
	\item цветовое кодирование модуля вектора.
\end{itemize}

В задачах вычислительной механики часто требуется визуализировать:

\begin{itemize}
	\item расчётную сетку;
	\item поле температуры;
	\item поле давления;
	\item скорость потока;
	\item напряжения и деформации;
	\item распределение ошибок;
	\item движение частиц;
	\item изменение решения во времени.
\end{itemize}

Большое значение имеет \textbf{интерактивная визуализация}. Пользователь может
изменять:

\begin{itemize}
	\item положение камеры;
	\item масштаб;
	\item параметры отображения;
	\item момент времени;
	\item отображаемую переменную;
	\item порог изоповерхности;
	\item диапазон цветовой шкалы.
\end{itemize}

Типичный интерактивный графический цикл имеет вид:

\begin{enumerate}
	\item обработать события ввода;
	\item изменить состояние сцены;
	\item выполнить необходимые вычисления;
	\item сформировать новый кадр;
	\item вывести его на экран.
\end{enumerate}

При интерактивной визуализации важна частота обновления изображения
(\textit{frame rate}). Если за секунду выводится $N$ кадров, то
\[
FPS=N.
\]

При частоте $60$ FPS на построение одного кадра приходится примерно
\[
\frac{1}{60}\approx 16.7\text{ мс}.
\]

Поэтому при работе с большими расчётными данными применяются:

\begin{itemize}
	\item уменьшение объёма отображаемых данных;
	\item уровни детализации;
	\item пространственные структуры данных;
	\item параллельная обработка;
	\item вычисления на GPU;
	\item предварительная обработка данных.
\end{itemize}

Таким образом, графические библиотеки отделяют прикладную программу от деталей
конкретного графического оборудования, а современный графический конвейер
обеспечивает массово-параллельное преобразование геометрии и данных в изображение.
Для задач математического моделирования компьютерная графика играет двойную роль:
она используется как средство построения пользовательского интерфейса и как
инструмент исследования результатов численного эксперимента.

\newpage